\section{Praktikum}
\label{sec:praktikum-bab7}

\begin{praktikumbox}
\vspace{1.1em}
\textbf{Judul:} Menulis, meng-assemble, men-link, dan menjalankan
program \file{.EXE} serta \file{.COM}.\\[0.3em]
\textbf{Sub-CPMK:} 2.1 \quad\textbullet\quad
\textbf{Alokasi waktu:} 100 menit \quad\textbullet\quad
\textbf{Bentuk:} individu
\end{praktikumbox}

\subsection{Tujuan Praktikum}

Setelah menyelesaikan praktikum ini, mahasiswa mampu menulis program
assembly, meng-assemble dengan TASM, men-link dengan TLINK,
menjalankan program \file{.EXE} dan \file{.COM}, serta memeriksa
berkas hasil dengan DEBUG.

\subsection{Alat dan Bahan}

\begin{itemize}
    \item Komputer dengan emu8086 terpasang.
    \item DOSBox dengan TASM, TLINK, dan DEBUG.
    \item Berkas \file{hello.asm} dan \file{direktif-data.asm} dari
          \file{code/asm/}.
    \item Lembar jurnal praktikum.
\end{itemize}

\subsection{Langkah Kerja}

\begin{enumerate}
    \item \textbf{Menyiapkan program \file{.EXE}.}
    \begin{enumerate}
        \item Buka DOSBox, lalu masuk ke direktori kerja.
        \item Salin \file{hello.asm} ke direktori kerja.
        \item Rakit dengan \file{tasm hello.asm}.
        \item Link dengan \file{tlink hello.obj}.
        \item Jalankan \file{hello} dan catat keluarannya.
    \end{enumerate}

    \item \textbf{Memeriksa berkas hasil.}
    \begin{enumerate}
        \item Catat ukuran \file{hello.obj} dan \file{hello.exe}.
        \item Jelaskan mengapa berkas \file{.OBJ} belum dapat
              dijalankan.
        \item Jelaskan peran TLINK dalam menghasilkan
              \file{hello.exe}.
    \end{enumerate}

    \item \textbf{Membuat program \file{.COM}.}
    \begin{enumerate}
        \item Tulis program berikut sebagai \file{hello-com.asm}.

\begin{lstlisting}[language=MASM, numbers=none]
.MODEL TINY
.CODE
ORG 100h
START:
    JMP MULAI
pesan DB 'Halo, COM!', 0Dh, 0Ah, '$'
MULAI:
    MOV AH, 09h
    MOV DX, OFFSET pesan
    INT 21h
    INT 20h
END START
\end{lstlisting}

        \item Rakit dengan \file{tasm hello-com.asm}.
        \item Link dengan \file{tlink /t hello-com.obj}.
        \item Jalankan \file{hello-com} dan catat keluarannya.
        \item Bandingkan ukuran \file{.COM} dengan \file{.EXE}.
    \end{enumerate}

    \item \textbf{Memeriksa program dengan DEBUG.}
    \begin{enumerate}
        \item Jalankan \file{debug hello.exe}.
        \item Ketik \file{r} dan catat nilai \reg{CS}, \reg{DS},
              dan \reg{IP}.
        \item Ketik \file{d ds:0000 L 20} untuk melihat isi segmen
              data.
        \item Temukan string \code{"Halo, dunia!"} pada hasil
              pemeriksaan.
        \item Keluar dari DEBUG dengan \file{q}.
    \end{enumerate}

    \item \textbf{Menguji direktif data.}
    \begin{enumerate}
        \item Rakit dan link \file{direktif-data.asm}.
        \item Jalankan program dan catat keluarannya.
        \item Buka program dengan DEBUG, lalu periksa isi segmen
              data dengan \file{d ds:0000 L 40}.
        \item Tandai byte yang berasal dari \dir{DB}, \dir{DW},
              \dir{DD}, dan \dir{DUP}.
    \end{enumerate}

    \item \textbf{Memodifikasi program.}
    \begin{enumerate}
        \item Ubah isi \code{pesan} pada \file{hello.asm} menjadi
              nama lengkap Anda.
        \item Ubah \code{bintang DB 4 DUP('*')} menjadi
              \code{6 DUP('*')}.
        \item Rakit ulang kedua program dan pastikan keluarannya
              berubah.
    \end{enumerate}
\end{enumerate}

\subsection{Hasil yang Dikumpulkan}

\begin{enumerate}
    \item Berkas \file{hello.asm}, \file{hello-com.asm}, dan
          \file{direktif-data.asm} versi modifikasi.
    \item Tangkapan layar keluaran ketiga program.
    \item Tabel perbandingan ukuran \file{.OBJ}, \file{.EXE}, dan
          \file{.COM}.
    \item Tangkapan layar pemeriksaan segmen data dengan DEBUG.
    \item Penjelasan singkat perbedaan program \file{.COM} dan
          \file{.EXE}.
\end{enumerate}

\begin{asesmenbox}
\vspace{1.1em}
\textbf{Asesmen praktikum 5 (bobot 5\% dari komponen praktikum).}

\begin{tabular}{@{}p{6.2cm}cp{4.6cm}@{}}
    \toprule
    \textbf{Indikator} & \textbf{Bobot} & \textbf{Kriteria} \\
    \midrule
    Program \file{.EXE} berhasil dirakit dan dijalankan & 30\% &
    Berkas \file{.OBJ} dan \file{.EXE} terbentuk, keluaran benar \\
    \addlinespace
    Program \file{.COM} berhasil dibuat & 25\% &
    Berkas \file{.COM} terbentuk dan berjalan dengan benar \\
    \addlinespace
    Pemeriksaan segmen data dengan DEBUG & 25\% &
    String dan data ditemukan serta ditafsirkan dengan benar \\
    \addlinespace
    Analisis perbedaan COM dan EXE & 20\% &
    Perbedaan segmen, ukuran, dan inisialisasi DS dijelaskan \\
    \bottomrule
\end{tabular}
\end{asesmenbox}
